% =============================================================================
% ch03_sensors.tex — Chapter 3: Sensor Modalities | NavigatorCrew
% Compiler: XeLaTeX | Language: Hebrew primary, English technical terms via \en{}
% =============================================================================

\selectlanguage{hebrew}

\section{מודאליות החיישנים}
\label{sec:sensors}

% ── 3.1 LiDAR Sensors ───────────────────────────────────────────────────────
\subsection{חיישני \en{LiDAR}: עקרונות פעולה ומאפיינים}
\label{subsec:lidar}

חיישני \en{LiDAR} (\en{Light Detection and Ranging}) מהווים את
עמוד השדרה של מערכות מיפוי וניווט תלת-ממדיות לרחפנים. עקרון
הפעולה מבוסס על שיגור פולסי לייזר ומדידת זמן החזרתם, בדומה
לעקרון האקולוקציה הביולוגית. חיישני \en{LiDAR} מודרניים,
כגון \en{Velodyne VLP-16} ו-\en{Ouster OS-0}, מספקים ענני
נקודות בקצב של 300,000 עד 2.2 מיליון נקודות לשנייה, עם טווח
פעולה של 50 עד 200 מטר.

שתי גישות עיקריות קיימות בתחום ה-\en{LiDAR} לרחפנים. הגישה
הראשונה היא \en{LiDAR} מכני, שבו סורק מסתובב פולט קרני לייזר
בזוויות שונות. גישה זו מספקת שדה ראייה של 360 מעלות אך סובלת
ממשקל גבוה ומחלקים מכניים רגישים. הגישה השנייה היא \en{LiDAR}
מבוסס \en{MEMS} (\en{Micro-Electro-Mechanical Systems}), שבו
מיקרו-מראה מתנודדת מכוונת את קרן הלייזר. חיישנים אלו קלים
וקטנים יותר, אך שדה הראייה שלהם מוגבל ל-90-120 מעלות.

\begin{figure}[H]
    \centering
    \includegraphics[width=0.92\columnwidth]{figures/fig_lidar_scanning_mechanisms.png}
    \caption{השוואה בין מנגנוני סריקת \en{LiDAR}: מכני (שמאל) לעומת
             \en{MEMS} (ימין). המנגנון המכני מספק 360 מעלות אך כבד יותר;
             ה-\en{MEMS} קל וקטן יותר אך בעל שדה ראייה מוגבל.}
    \label{fig:lidar_scanning_mechanisms}
\end{figure}

% ── 3.2 IMU Sensors ─────────────────────────────────────────────────────────
\subsection{חיישני \en{IMU}: מדידת תאוצה ומהירות זוויתית}
\label{subsec:imu}

יחידת המדידה האינרציאלית (\en{IMU}) מספקת מדידות תאוצה ומהירות
זוויתית בקצב גבוה, הנע בין 100 ל-1000 הרץ. \en{IMU} טיפוסי
לרחפנים כולל מד תאוצה תלת-צירי וגירוסקופ תלת-צירי, המבוססים
על טכנולוגיית \en{MEMS}. דיוק המדידות תלוי באיכות החיישן:
\en{IMU} תעשייתי (כגון \en{ADIS16470}) מציג יציבות אפס של
0.1 מעלות לשנייה, בעוד \en{IMU} זול (כגון \en{MPU9250})
סובל מסחיפה של 1-5 מעלות לשנייה.

האתגר המרכזי בשימוש ב-\en{IMU} הוא הצטברות שגיאות לאורך זמן
(\en{Drift}). שגיאות אלו נובעות מרעש לבן במדידות התאוצה,
משגיאת אפס (\en{Bias}) המשתנה עם הטמפרטורה, ומרעש דגימה
(\en{Random Walk}). כדי להתמודד עם סחיפה זו, יש צורך בשילוב
ה-\en{IMU} עם חיישנים מוחלטים כגון \en{GPS} או \en{LiDAR},
בתהליך המכונה \en{Sensor Fusion} \cite{Thrun2005ProbRobotics}.

\begin{table}[H]
\centering
\begin{tabular}{lccc}
\toprule
סוג \en{IMU} & יציבות אפס & רעש תאוצה & עלות \\
\midrule
תעשייתי (\en{ADIS16470}) & 0.1$^\circ$/ש׳ & 0.01 מ/ש$^2$/√הרץ & 500-1000 \$ \\
בינוני (\en{ICM-20948}) & 0.5$^\circ$/ש׳ & 0.05 מ/ש$^2$/√הרץ & 50-100 \$ \\
זול (\en{MPU9250}) & 1-5$^\circ$/ש׳ & 0.1 מ/ש$^2$/√הרץ & 5-15 \$ \\
\bottomrule
\end{tabular}
\caption{השוואת מאפייני \en{IMU} נפוצים לרחפנים.}
\label{tab:imu_characteristics}
\end{table}

% ── 3.3 Camera Sensors ──────────────────────────────────────────────────────
\subsection{חיישני מצלמה: עיבוד תמונה וזיהוי תכונות}
\label{subsec:cameras}

מצלמות \en{RGB} ומצלמות עומק (\en{RGB-D}) מספקות מידע חזותי
עשיר החיוני לזיהוי מכשולים, לוקליזציה ומיפוי. בניגוד ל-\en{LiDAR},
המספק מידע תלת-ממדי מדויק אך דל מבחינה סמנטית, מצלמות מאפשרות
זיהוי של עצמים, טקסטורות ומרקמים. חיישני מצלמה נפוצים לרחפנים
כוללים את \en{Intel RealSense D435} ו-\en{Flir Blackfly S},
המספקים רזולוציה של 1-20 מגה-פיקסל בקצב של 30-90 פריימים לשנייה.

שיטת \en{ORB-SLAM3} \cite{MurArtal2015ORB} מדגימה כיצד ניתן
להשתמש בתכונות \en{ORB} (\en{Oriented FAST and Rotated BRIEF})
למיפוי ולוקליזציה בזמן אמת. תכונות אלו עמידות לשינויים בסיבוב,
בקנה מידה ובתאורה, ומאפשרות התאמה בין תמונות עוקבות גם בתנאי
תנועה מהירה. עם זאת, מערכות חזותיות סובלות מכשלים בסביבות דלות
טקסטורה, בתנאי תאורה משתנים ובתנועה סיבובית טהורה.

\begin{equation}
    \mathbf{f}_{\text{ORB}} = \{ \mathbf{p}_i, \mathbf{d}_i \}_{i=1}^{N},
    \quad \mathbf{p}_i \in \mathbb{R}^2,
    \quad \mathbf{d}_i \in \{0,1\}^{256}
    \label{eq:orb_features}
\end{equation}

כאשר $\mathbf{p}_i$ הוא מיקום התכונה בתמונה ו-$\mathbf{d}_i$ הוא
וקטור התיאור הבינארי באורך 256 סיביות. התאמה בין תכונות בשתי
תמונות עוקבות מתבצעת באמצעות מרחק האמינג (\en{Hamming distance}).

% ── 3.4 Multi-Modal Fusion ──────────────────────────────────────────────────
\subsection{שילוב רב-מודאלי: אתגרים והזדמנויות}
\label{subsec:multimodal}

שילוב חיישנים מרובים (\en{Multi-Sensor Fusion}) מאפשר ניצול
היתרונות של כל מודאליות תוך כיסוי חולשותיה. הגישה המקובלת היא
\en{Fusion} הדוק (\en{Tightly-Coupled Fusion}), שבו מדידות
גולמיות מכל החיישנים מאופטימיזות במשותף בפונקציית עלות אחת.
\en{LIO-SAM} \cite{Thrun2005ProbRobotics} מדגים גישה זו על ידי
שילוב \en{LiDAR}, \en{IMU} ו-\en{GPS} במסגרת \en{Factor Graph}.

האתגר המרכזי בשילוב רב-מודאלי הוא כיול מרחבי וזמני מדויק בין
החיישנים. שגיאת כיול קטנה של 0.1 מעלות בסיבוב או 1 ס״מ בתרגום
עלולה לגרום לסחיפה משמעותית לאורך מסלולים ארוכים. אתגר נוסף הוא
טיפול בנשירת חיישנים (\en{Sensor Dropout}): כאשר \en{GPS} אובד
בקניונים עירוניים או \en{LiDAR} נכשל בערפל, המערכת חייבת
להתדרדר באופן הדרגתי ולא להתבדר.

\begin{figure}[H]
    \centering
    \includegraphics[width=0.92\columnwidth]{figures/fig_multi_sensor_fusion_architecture.png}
    \caption{ארכיטקטורת \en{Fusion} רב-מודאלי במערכת \en{NavigatorCrew}.
             חיישני \en{LiDAR}, \en{IMU}, מצלמה ו-\en{GPS} מזינים
             במקביל את מסגרת ה-\en{Factor Graph} לאופטימיזציה משותפת.}
    \label{fig:multi_sensor_fusion}
\end{figure}

\begin{equation}
    \mathbf{x}^* = \arg\min_{\mathbf{x}} \sum_{i} \|e_i(\mathbf{x})\|_{\Omega_i}^2
    + \sum_{j} \|e_j(\mathbf{x})\|_{\Omega_j}^2
    + \sum_{k} \|e_k(\mathbf{x})\|_{\Omega_k}^2
    \label{eq:multimodal_cost}
\end{equation}

משוואה~\ref{eq:multimodal_cost} מציגה את פונקציית העלות המשותפת,
כאשר האינדקסים $i$, $j$, $k$ מייצגים אילוצי \en{LiDAR},
\en{IMU} ומצלמה בהתאמה, ו-$\Omega$ הן מטריצות האינפורמציה
המתאימות.

% ── 3.5 Summary ─────────────────────────────────────────────────────────────
\subsection{סיכום}
\label{subsec:sensors_summary}

בפרק זה סקרנו את שלוש מודאליות החיישנים המרכזיות במערכת
\en{NavigatorCrew}: \en{LiDAR} למיפוי תלת-ממדי מדויק,
\en{IMU} לאומדן תנועה בתדר גבוה, ומצלמות למידע חזותי עשיר.
השילוב ההדוק בין מודאליות אלו, כפי שיוצג בפרקים הבאים,
מאפשר התגברות על מגבלות כל חיישן בנפרד.